\section{Validation of floating-point bounds}\label{sec:numerical-validation}

This appendix gives the row correction and dual-bound evaluation used by the
numerical propagator of Section~\ref{sec:algorithms-policy}. Each stored
floating-point number is treated as the exact dyadic rational it represents.
The LP variables have a finite box of bounds.

\begin{proposition}[Outward row correction]\label{prop:row-correction}
Let $a^\top z\le\beta$ be valid for a set $S$ contained in the box
$Z=[\underline z,\overline z]$. For any vector $\tilde a$, every $z\in S$
satisfies $\tilde a^\top z\le\tilde\beta$ whenever
\begin{equation}\label{eq:row-correction}
 \tilde\beta\ \ge\ \beta+\sum_j
 \max\bigl\{(\tilde a_j-a_j)\underline z_j,\,
            (\tilde a_j-a_j)\overline z_j\bigr\}.
\end{equation}
\end{proposition}

\begin{proof}
For $z\in S$,
$\tilde a^\top z=a^\top z+\sum_j(\tilde a_j-a_j)z_j\le\beta+
\sum_j\max_{\underline z_j\le s\le\overline z_j}(\tilde a_j-a_j)s$, and a linear
function of $s$ attains its maximum over an interval at an endpoint.
\end{proof}

The implementation takes $\tilde a_j$ to be the floating-point number nearest
to $a_j$, evaluates the right-hand side of~\eqref{eq:row-correction} in exact
rational arithmetic over the LP's variable bounds, and rounds it upward to a
floating-point $\tilde\beta$. The coefficients of the original rows, of the
McCormick rows, and of the tangents are floating-point numbers already, while
the secant slope $\ell_i+u_i$ may need correction; right-hand sides such as
$\ell_i\ell_j$ and $p^2$ stay exact until the final upward rounding.
Proposition~\ref{prop:row-correction} applies to each product row with $S$
the set of exact lifts $(x,(x_ix_j))$ of the points $x\in B$, and to each
original row and the cutoff row with $S$ the set of exact lifts of the
points of $B$ that satisfy that row. Both sets lie in the LP's variable box
by construction.

The LP solver returns approximate inequality multipliers. They need not be
dual feasible. The following bound uses any nonpositive multipliers and absorbs
the dual residual through the variable box. Bounds of this kind, evaluated
with directed rounding, are an established way to obtain safe LP bounds
from approximate dual solutions~\cite{neumaier2004-safe-bounds-in-linear-and};
we state the form we use. It is the
minimization form of
Proposition~\ref{prop:corrected-dual-con}, which applies the same idea to a
stored dual vector after the rows have changed; here the rows are current and
the multipliers inexact.

\begin{proposition}[Dual bound with a residual term]\label{prop:dual-residual}
Consider $\mu=\min\{d^\top z:Az\le b,\ z\in Z\}$ with a finite box
$Z=[\underline z,\overline z]$, and $\mu=+\infty$ if the LP is infeasible.
For every $y\le0$, with $r=d-A^\top y$,
\begin{equation}\label{eq:dual-residual}
 \mu\ \ge\ y^\top b+\sum_j\min\{r_j\underline z_j,\,r_j\overline z_j\}.
\end{equation}
The right-hand side is the LP dual function at $y$, with the multipliers of
the variable bounds eliminated; its maximum over $y\le0$ equals $\mu$ when the
LP is feasible.
\end{proposition}

\begin{proof}
If the LP is infeasible, there is nothing to prove.
For feasible $z$, $y\le0$ and $Az\le b$ give $y^\top Az\ge y^\top b$, so
$d^\top z=y^\top Az+r^\top z\ge y^\top b+r^\top z$. Minimizing $r^\top z$
separately over each interval gives~\eqref{eq:dual-residual}. For the second
statement, the LP dual of $\mu$ has multipliers $y\le0$ for $Az\le b$ and
$\alpha,\gamma\ge0$ for $z\ge\underline z$ and $-z\ge-\overline z$, with
$d=A^\top y+\alpha-\gamma$ and objective
$y^\top b+\underline z^\top\alpha-\overline z^\top\gamma$. For fixed $y$,
the best choice is $\alpha_j=\max\{r_j,0\}$ and $\gamma_j=\max\{-r_j,0\}$, which
gives the right-hand side of~\eqref{eq:dual-residual} because
$\underline z_j\le\overline z_j$. LP strong duality completes the proof.
\end{proof}

The implementation clips the returned multipliers to $y=\min\{\tilde y,0\}$ and
evaluates~\eqref{eq:dual-residual} with directed rounding; it rejects
multipliers that are not finite. Because the proposition holds for every
$y\le0$, the validity of the result depends neither on the accuracy of the
multipliers nor on the solver's sign convention for them: poor or wrongly
signed multipliers only weaken the bound.

Let $\bar F$ be the set of binary64 numbers together with $\pm\infty$, and
let $\Omega$ be its largest finite element. For $t\in\bar F$, let
$\operatorname{pred}(t)$ and $\operatorname{succ}(t)$ be the next smaller and
the next larger element of $\bar F$, the IEEE~754 operations nextDown and
nextUp, with $\operatorname{pred}(-\infty)=-\infty$ and
$\operatorname{succ}(+\infty)=+\infty$. Thus
$\operatorname{pred}(+\infty)=\Omega$ and $\operatorname{succ}(-\infty)=-\Omega$,
and both maps are nondecreasing. For an addition, subtraction, or
multiplication of elements of $\bar F$ whose exact result $s$ is defined in
the extended reals, write $\operatorname{fl}(s)$ for the computed result. We
assume that rounding is \emph{faithful}: $\operatorname{fl}(s)=s$ if
$s\in\bar F$, and otherwise $\operatorname{fl}(s)$ is one of the two
consecutive elements of $\bar F$ that enclose $s$.

\begin{lemma}[One-step enclosure]\label{lem:enclosure}
Under faithful rounding,
$\operatorname{pred}(\operatorname{fl}(s))\le s\le
\operatorname{succ}(\operatorname{fl}(s))$ for every exact result $s$.
\end{lemma}

\begin{proof}
If $\operatorname{fl}(s)=s$, the claim holds because
$\operatorname{pred}(t)\le t\le\operatorname{succ}(t)$ for every $t\in\bar F$.
Otherwise let $s^-<s<s^+$ be the consecutive elements of $\bar F$ that
enclose $s$. If $\operatorname{fl}(s)=s^-$, then
$\operatorname{pred}(s^-)\le s^-<s<s^+=\operatorname{succ}(s^-)$; the case
$\operatorname{fl}(s)=s^+$ is symmetric.
\end{proof}

With gradual underflow, every IEEE~754 rounding direction is faithful,
including overflow to $\pm\infty$ or to $\pm\Omega$. The argument therefore
also covers a host solver that leaves a directed rounding mode active.
Flush-to-zero and denormals-are-zero modes are not faithful, and the
argument assumes that they are disabled.

The evaluation computes three enclosures, each preserved by
Lemma~\ref{lem:enclosure}; the first two are accumulated in one sweep over
the rows with $y_r\ne0$.
\begin{enumerate}
\item A lower bound $\sigma$ of $y^\top b$: starting from $\sigma=0$, add
$\operatorname{pred}(\operatorname{fl}(y_rb_r))$ for every row with
$y_r\ne0$ and apply $\operatorname{pred}$ to every rounded sum.
\item An interval $[\underline r_j,\overline r_j]$ containing $r_j$: start
from $d_j$, and for each nonzero $a_{rj}$ with $y_r\ne0$, let
$t=\operatorname{fl}(a_{rj}y_r)$ and update
$\underline r_j\leftarrow\operatorname{pred}(\operatorname{fl}(\underline r_j-\operatorname{succ}(t)))$
and
$\overline r_j\leftarrow\operatorname{succ}(\operatorname{fl}(\overline r_j-\operatorname{pred}(t)))$.
\item The bound $\tilde\mu$: starting from $\sigma$, add, for each $j$,
$\operatorname{pred}$ of the smallest of the four rounded products of
$\{\underline r_j,\overline r_j\}$ and $\{\underline z_j,\overline z_j\}$,
again applying $\operatorname{pred}$ to every rounded sum.
\end{enumerate}
Overflow cannot produce an undefined operation in the first two steps. The
lower bounds $\sigma$ and $\underline r_j$ start finite and are afterwards
results of $\operatorname{pred}$, so they are never $+\infty$, and neither
is the term $\operatorname{pred}(\operatorname{fl}(y_rb_r))$ added to
$\sigma$; the term $\operatorname{succ}(t)$ subtracted from $\underline r_j$
is never $-\infty$. Symmetrically, $\overline r_j$ is never $-\infty$, and
$\operatorname{pred}(t)$ is never $+\infty$. No sum or difference therefore
has the form $\infty-\infty$. An infinite intermediate
value is $-\infty$ in a lower bound or $+\infty$ in an upper bound, which
remains valid. The third step rejects the computation if any of its products
or the final value is not finite; an accepted result thus uses finite
$\underline r_j$ and $\overline r_j$. For each $j$, let $p_j$ be the smallest
of the four exact vertex products. The smallest rounded product is at most
$\operatorname{fl}(p_j)$, and $\operatorname{pred}$ is nondecreasing, so the
term added for $j$ is at most $\operatorname{pred}(\operatorname{fl}(p_j))\le p_j$.
Since the bilinear function $\rho s$ attains its minimum over a rectangle at
a vertex, an accepted result satisfies
\[
 \tilde\mu\ \le\ y^\top b+\sum_j
 \min_{\rho\in[\underline r_j,\overline r_j],\,s\in[\underline z_j,\overline z_j]}\rho s
 \ \le\ y^\top b+\sum_j\min\{r_j\underline z_j,r_j\overline z_j\}\ \le\ \mu .
\]
These arguments assume IEEE binary64 arithmetic with gradual underflow.
They establish a bound for the stored floating-point LP.
Together with Proposition~\ref{prop:row-correction}, that LP is a relaxation of
the exact product graph of the stored model on $B$. They do not relate the
stored binary model to a decimal or symbolic model upstream of its import.
